% 本文件由 LaTeX 排版 Agent 根据有效 translation.json 自动生成。
% author=Peter of Alexandria; work_id=0620; title=正典书信

\chapter{正典书信}

% structure: p0001 source=h1 target=omitted reason=page-title-already-rendered-by-parent

\Emph{附\Person{狄奥多尔·巴尔萨蒙}与\Person{若望·佐纳拉斯}的注释。}\par

\WorkTitle{真福伯多禄，亚历山大里亚总主教，在其忏悔讲道中所载的教规}。\par

% structure: p0004 source=h2 target=section reason=relative-heading-rank; base=h2

\section{第一教规}

然而，既然迫害的第四个逾越节已经来到，对于那些曾被逮捕、投入监牢，忍受了难以承受的酷刑、不可忍受的鞭打以及其他许多可怕的折磨，后来因肉身的软弱而背叛的人——即使他们起初因其严重的跌倒而不被接纳——但因他们曾奋力抵抗并长久坚持（他们并非自愿如此，而是因肉身的软弱而背叛，因为他们在自己身上带着耶稣的印记），并且有些人如今已为他们的过错哀哭了三年：我认为，从他们谦卑地前来开始，应再补加四十天，使他们记念这些事；这四十天期间，我们的主和救主耶稣基督虽然禁食，却在受洗后被魔鬼试探。当他们在这期间刻苦操练、恒切禁食之后，就让他们警醒祈祷，默想主对那试探祂、叫祂俯伏朝拜者所说的话：\BibleQuote{撒殚，退到我后面去！因为经上记载：你要朝拜主，你的天主，惟独事奉祂。} \BibleRef{玛 4:10}\par

\Person{巴尔萨蒙}。本教规论及在迫害中否认信仰并正在做补赎的人。第一教规规定：那些在遭受许多折磨后因软弱未能坚持而祭拜偶像、并在忏悔中度过了三年的人，应再加四十天，然后接纳他们进入教会。请注意这些教规，它们为那些否认天主并寻求悔改的人，以及那些自愿殉道、跌倒后又重新宣认信仰的人，还有其他类似情况，制定了各种有益的规则。也请查阅安西拉会议的许多教规，因为这样有益于你们。\par

\Person{佐纳拉斯}。圣伯多禄在那些动乱时期否认信仰的人中作出区分，并说：对那些曾被带到暴君面前、投入监牢、忍受了极其严重的酷刑和不可忍受的鞭打，以至于任何照料或药物都无法治愈（\OriginalTerm{因为}{\GreekText{ἀνη}}\par

\SourceRecord{Translated by James B.H. Hawkins.；Ante-Nicene Fathers；Vol. 6.；Edited by Alexander Roberts, James Donaldson, and A. Cleveland Coxe.；Buffalo, NY: Christian Literature Publishing Co.,；1886.；<http://www.newadvent.org/fathers/0620.htm>.}
